
\section{Three-prime exclusion}\label{sec:threeprime}

\begin{proof}[Proof of Theorem~\ref{main:threeprime}]
Divide the roots of a positive square progression by their gcd $g$.
The primitive roots $r',s',t'$ are all odd (reduce modulo $4$), and
$U=(t'-r')/2$, $V=(t'+r')/2$, $s'$ form a primitive right triangle:
$U^2+V^2=s'^2$, $\gcd(U,V)=1$. Its area $A=UV/2$ satisfies
$d/4=Ag^2$, so its prime support has at most three elements.

Aebi's theorem and closing remark~\cite{AEB26} classify all such
primitive triangles, up to exchanging the legs. Their areas and
squarefree parts are, respectively,
\[
\begin{array}{c|rrrrrrr}
A&6&30&60&180&84&504&1224\\\hline
\operatorname{sf}(A)&6&30&15&5&21&14&34
\end{array}
\]
and each area has exactly one primitive triangle. The squarefree parts
are distinct. If two positive progressions had the same $d$, the ratio
of their primitive areas would be a rational square, forcing the same
area and the same $g$. The triangle then recovers the unique roots
$g|U-V|,gs',g(U+V)$, so the progressions coincide. This proves the
statement and its array consequences. The bound on the number of primes
cannot be raised to four for this uniqueness statement: at $d=840$ the
two triples $(1,29,41)$ and $(23,37,47)$ are distinct examples.
\end{proof}

\subsection{The elementary endpoint route}

For comparison, we retain the local endpoint argument and valuation
restrictions. In the rest of this section the configuration is primitive
and satisfies the full-magic interaction.

Throughout, $\supp(2d)=\{2,3,p\}$ with $p\neq3$ prime, and
$d/2=XYZ$ with $X=2^{\alpha}$, $Y=3^{\beta}$, $Z=p^{\gamma}$. The two
technical stages are stated first; the headline
Theorem~\ref{main:threeprime} subsumes both by a different global
argument, and its proof is independent of theirs.

\begin{proposition}[All-endpoint case, $p\equiv3\bmod8$]\label{prop:threeblock}
There is no primitive positive Parker configuration with
$\supp(2d)=\{2,3,p\}$ and $p\equiv3\pmod8$.
\end{proposition}

\begin{proof}
By the endpoint table of Section~\ref{sec:supportlaw}, $p\equiv3\bmod8$
forces all three indices to be endpoints at $p$, at $3$ (same residue
argument), and --- by the opposite-parity mechanism at $2$ --- at the
$2$-power as well. Each progression's equal-area Pythagorean triangle is
then primitive, and its legs are a unitary partition of the three whole
blocks $X,Y,Z$. Apart from the impossible split $\{1,XYZ\}$, only
$\{X,YZ\}$, $\{Y,XZ\}$, $\{Z,XY\}$ remain, and three pairwise-distinct
progressions must use all three. Factoring the hypotenuse equations of the
last two gives
\[
Y^2-1=2XZ,\qquad Z^2-1=2XY ,
\]
whose difference factors as $(Y-Z)(Y+Z+2X)=0$. Positivity forces $Y=Z$,
after which $Y(Y-2X)=1$, impossible for positive prime-power blocks.
\end{proof}

\begin{proposition}[One-interior law, $p\equiv5,7\bmod8$]\label{prop:oneinterior}
Let $\supp(2d)=\{2,3,p\}$ with $p\equiv5$ or $7\pmod8$. The all-endpoint
branch is impossible by the three-block argument, so $\gamma=1$ is
excluded, and for $\gamma\ge2$ exactly one index is interior at $p$. If
$p\equiv5\pmod8$ the centre progression is the interior one and exactly one
of
\[
v_p(X^2+Y^2)=2h,\qquad v_p(1+Y^2)=2h,\qquad v_p(1+X^2)=2h
\]
holds, with $h=v_p(s_0)\ge1$; if $p\equiv7\pmod8$ one outer progression is
interior and exactly one of the six ordered differences
\[
v_p(2X^2-Y^2),\; v_p(2Y^2-X^2),\; v_p(2-Y^2),\;
v_p(2-X^2),\; v_p(2Y^2-1),\; v_p(2X^2-1)
\]
equals $2h$, with $h=\min(a,\gamma-a)$. These are necessary branch laws,
not contradictions; their exact certificate covers all nine identities.
\end{proposition}

The nine residual branches of Proposition~\ref{prop:oneinterior} need not
be analyzed one by one: the headline theorem closes every case, including
$p\equiv1\bmod8$, at once.

\begin{proof}[An elementary proof for primitive exact support $\{2,3,p\}$]
By the endpoint table, every index is an endpoint at $3$ and at least two
indices are endpoints at $p$. Those two indices give two \emph{distinct
primitive} Pythagorean triangles of the same area $d/4$, whose area has
exactly three distinct prime factors. Writing $X=2^{\alpha}=2x$, the three
endpoint types satisfy
\[
YZ=x^2-1,\qquad Y^2-1=4xZ,\qquad Z^2-1=4xY .
\]
Combining the first with each of the others reduces both mixed pairs to
one equation
\[
W(W^2-1)=4x(x^2-1),
\]
where $W$ is the remaining odd block. The $2$-adic valuation of the right
side forces $W\equiv\pm1\pmod{2x}$, while monotonicity of $W(W^2-1)$ gives
$x<W<2x$; the only residue in that window is $W=2x-1$, which forces
$x=2$, $W=3$, and then the third block equals $1$ --- not a prime power in
the support. The remaining pair of endpoint types forces an equality of
powers of two distinct primes, equally impossible. Hence no primitive
positive Parker configuration has $\supp(2d)=\{2,3,p\}$ for any prime
$p\neq3$.
\end{proof}

\begin{remark}[Attribution and scope]
The headline theorem is the squareclass consequence of Aebi's complete
triangle classification proved at the start of this section. The endpoint
argument gives another proof of its primitive full-magic consequence,
without a priority claim. The uniqueness assertion fails already at four
primes, as the difference-$840$ example shows. Section~\ref{sec:smoothsupport}
nevertheless closes particular four-prime supports by a separate complete
catalogue; arbitrary four-prime supports and global support-size control
remain open.
\end{remark}
